\section{Verification Evidence}
\label{sec:appendix_verification_evidence}

This appendix records the verification evidence for the final writing pack. These outputs confirm that the reporting layer can be regenerated from existing evidence files; they are not new model-evaluation results.

\begin{itemize}
    \item Test command: \texttt{python -m pytest tests/prototype5 -v}. Key output: \texttt{============================= 62 passed in 0.97s ==============================}.
    \item Orchestrator command: \texttt{python -m src.prototype5.run_orchestrator}. Key output: \texttt{Prototype 5 Final Evidence Orchestrator: COMPLETE}.
    \item Final pack generator command: \texttt{python -m src.prototype5.generate_final_dissertation_pack}. Key output: \texttt{Prototype 5 final dissertation pack generated.}.
\end{itemize}

The raw verification snapshots are stored at:
\begin{itemize}
    \item \texttt{results/prototype5/verification\_snapshots/pytest\_prototype5\_output.txt}
    \item \texttt{results/prototype5/verification\_snapshots/orchestrator\_output.txt}
    \item \texttt{results/prototype5/verification\_snapshots/final\_pack\_generator\_output.txt}
    \item \texttt{results/prototype5/verification\_snapshots/verification\_summary.md}
\end{itemize}

The generated evidence files include the final evidence dashboard, claims matrix, dissertation metrics, evidence manifest, Microsoft brief alignment matrix, safety-latency frontier and final pack completion report. Remaining caveats include the absence of production robot safety validation, no statistical generalisation beyond the benchmark, no quantitative cloud cost accounting, no detected GPU/NPU counters in Mode D, and Prototype 1 being context-only evidence.
